\section{Planning across turns and responding to feedback}
\label{sec:temporal}

The temporal sample contains three decisions from each of eighty player-games, selected by a salted hash from an earlier pool of public-message-unflagged records. One LLM-assisted reader examined the complete applied public and private fields, excluding provider explanations. Of the 240 decisions, fifty-nine contain forward-looking own-action statements, thirty-six contain prospective opponent-capacity or outcome statements, and twenty-five refer to correction or feedback, with overlap. Grok accounts for thirty-nine of the fifty-nine own-action records and twenty-two of the thirty-six opponent records. Fifty-three private fields are absent, all for Opus. These counts characterize the selected records rather than estimate a corpus-wide planning rate.

Most statements are contingent or open-ended rather than scoreable commitments. Eight specify buying a named property on landing there. Six never meet a later matching purchase opportunity. One is carried out. The other, Grok's G16 Virginia plan, encounters an affordable purchase menu at d159 with \$371 against a \$160 price. Its note still favors purchase and its public message claims the pink set is closed, but the received call starts an auction. Gemini buys the deed for \$380. This is an observed failure of follow-through with conflicting channels, not an identified change of mind or memory failure.

A positive longer target appears in G9. Sol's turn-29 note plans three houses on each dark-blue deed as funds arrive. Construction and financing over subsequent turns reach that target on turn 60. The sequence demonstrates an expressed target and its later realization; it does not show that retaining the note caused the outcome. The G15 offers likewise respond to an opponent's stated objection, while other notes make checkable deductions about publicly visible cash and collateral. These are concrete forms of opponent attention, rather than calibrated predictions of hidden beliefs.

The complete uneven-build recurrence ledger follows 114 error decisions to 109 later build submissions. Of fifty successors after a decision containing a corrective request with the validator message, nine repeat the error. Of fifty-nine successors without that feedback, fifty-four repeat it; all fifty-nine belong to Opus, whose refusal path consumed the corrective opportunity. The contrast documents repair and recurrence under different response conditions. Because feedback and configuration coincide, it does not identify feedback-caused learning.

\section{Discussion}
\label{sec:discussion}

\subsection{What the participants negotiate}

Participants bargain over positions extending beyond the exchanged deed. They discuss a rival's construction reserves, the value of withholding a completing property and whether funding a distressed seller benefits a third-party creditor. The offer census places those considerations in context: outright refusal is more common when only the offer author completes a group, even within opening and later offers, while counters commonly preserve the exchange and ask for better terms. The transcripts show what that bargaining involves. In G15, a recipient refuses its own completion after the proposer addresses a liquidity objection. In G2, Sol quotes different prices to different buyers, while Opus advertises a lower ceiling than its private contingency permits. These are observable responses to interdependent economic stakes.

The conduct spans successful exchange, refusal, factual correction and changed assurances. G7's bluff settles at the advertised price. G19's recipient corrects a false claim and still accepts the exchange. G13's seller selectively presents its financing options while the buyer discusses the estate at risk. G18's buyer withdraws a repeated price and resumes bidding after the seller deteriorates. Public deception is permitted in this task, so these incidents describe strategic conduct under that contract. Their context shows both what the tactic sought and what happened next, without attributing the resulting price or victory to the tactic alone.

\subsection{From ownership to income and control}

The next question is whether a negotiated position becomes productive. The systems complete similar numbers of groups by trade---seventeen, eighteen, fifteen and thirteen---but build on sixteen, seventeen, seven and three in the first management window. Those early-built groups collect \$53,810 of \$57,378 in post-completion rent. Owners sometimes preserve cash or defer redemption; elsewhere the response path prevents a requested build. Following acquisition into financing and construction thus connects bargaining with the income the position actually produces.

Financial access also determines what can happen at a debt endpoint. Sol's G20 sequence realizes \$1,365 from sales and mortgages to meet a \$1,400 bill, whereas Gemini's G16 calculation correctly recognizes that all available own resources fall short. Across the sixty eliminations, fifty-six debts exceed even the generous asset bound. These endpoints establish the immediate financing position, leaving the quality of earlier spending as a separate question.

Inheritance then extends the economic consequences beyond the eliminated player. Eventual winners receive thirty-five of forty nonterminal estates, and thirty-eight of those forty eliminations follow rent on developed property. A productive group can therefore be followed by acquisition of the debtor's remaining assets. This pathway explains why participants discuss who would inherit a blocker or a complete group. G13's insurance pitch and G18's withdrawn offer make that concern explicit; G15's bank-purchased oranges and subsequent estates show the corresponding sequence in play. Bargaining, development and creditor succession belong in the same account because the assets and incentives carry through all three stages.

\subsection{Scope of the motivating research questions}

Table~\ref{tab:questions} distinguishes direct population descriptions, contextual observations and constructs that remain unmeasured. An unresolved construct is not a negative capability result.

\begingroup\footnotesize
\renewcommand{\arraystretch}{1.09}
\begin{longtable}{@{}L{0.22\linewidth}L{0.36\linewidth}L{0.36\linewidth}@{}}
\caption{The eight motivating questions and three broader questions receive different kinds of evidence. The table is a scope map, not a benchmark scorecard.}\label{tab:questions}\\
\toprule
Question & Evidence in this cohort & Not established\\
\midrule\endfirsthead
\toprule
Question & Evidence in this cohort & Not established\\
\midrule\endhead
\bottomrule\endfoot
1. Strategic reasoning & Offer responses, stated rival consequences, local financing alternatives & Optimal actions, regret or exploitability\\
2. Opponent modeling & Public-capacity deductions and responses to stated objections & Calibrated hidden-belief or preference inference\\
3. Negotiation, cooperation and deception & Complete response census; selected bluffs, omission, correction, withdrawal and coordination appeal & Persuasion effects, fairness, joint surplus, agreement bias or deception prevalence\\
4. Long-horizon planning and memory & Acquisition-to-income trajectories; targets with later opportunities; documented memory conditions & Memory benefit or a general plan-completion rate\\
5. Adaptability and learning & Revisions, successful repairs and repeated errors under observed feedback & Feedback-caused learning or adaptation to unfamiliar opponents\\
6. Economics and risk & Realized flows, construction feasibility and all current-debt endpoints & Ex ante asset values, optimal reserves or risk-adjusted returns\\
7. Collective dynamics & Creditor concentration and a recognized financial-denial appeal & Coalition incidence, stability or a causal coordination effect\\
8. Validity and generalization & Reconstruction under the preserved implementation; explicit response and information conditions & A clean model ranking, corrected-interface performance or external transfer\\
\midrule
A. Strategic talk and action & Cases of realized targets, factual correction and conflicting execution & A general relation between explanation quality and strategic quality\\
B. General versus multi-agent capability & Not evaluated & A capability curve or transfer relationship\\
C. Skill, chance and opponent weakness & Descriptive seed/seat patterns and specific economic pathways & Their causal shares in victory or a counterfactual winner\\
\end{longtable}
\endgroup

Outcome patterns illustrate why the distinction matters. Sol wins one four-rotation block from every seat. In another block, all four configurations win from seat 2, whose player purchases both dark blues by turn 29. In G13--15, seat 1 completes orange through direct purchase on turn 47 under Opus, Sol and Gemini respectively, and wins each game. In the fourth rotation, G16, seat 1 is Grok: it completes orange by trade on turn 116 and loses. Thus seven games share early same-seat purchase paths across different configurations, while four follow Sol across seats. The purchases, subsequent decisions and random events are all part of the observed paths; these patterns do not establish that early dice decided seven outcomes. Across the cohort Sol wins eight games, Gemini seven, Opus three and Grok two. Neither pattern isolates strategic skill or chance, and no independent general-capability axis is measured.
\FloatBarrier
